\chapter{Overview}
\label{ch:overview}

Veilbox is a minimal, bootable operating system built entirely from source
using a custom-compiled Linux kernel with an embedded initramfs (initial RAM
filesystem). It boots directly to a BusyBox shell with containerd, runc,
nerdctl, and Dropbear SSH running out of the box. The entire operating system
is packaged as a single disk image --- 8 gigabytes raw or approximately 95
megabytes as a sparse VDI --- that can be imported into any virtual machine
manager or written directly to bare metal hardware.

\section{The Problem Veilbox Solves}

Building a Linux system from scratch has traditionally been a complex,
time-consuming process requiring deep expertise in kernel compilation,
initramfs generation, bootloader configuration, and userspace assembly.
Most solutions fall into three categories:

\begin{enumerate}[nosep]
  \item \textbf{General-purpose distributions} (Fedora, Ubuntu, Debian):
        Large, complex, with hundreds of packages and services. Installing
        a minimal container host requires post-install configuration, package
        removal, and security hardening.
  \item \textbf{Embedded build systems} (Buildroot, Yocto): Powerful but
        complex. Buildroot requires learning its Kconfig-based configuration
        system. Yocto requires understanding BitBake recipes, layers, and
        metadata. Both produce excellent results but have steep learning
        curves and long initial build times.
  \item \textbf{Minimal distributions} (Alpine Linux): Lightweight but still
        include a package manager, init system, and package repositories.
        Alpine uses OpenRC init and BusyBox, but still requires post-install
        configuration and has a 50+ MB root filesystem.
\end{enumerate}

Veilbox takes a different approach: it is a single-purpose operating system
designed for one task --- running containers in a virtual machine --- with no
extraneous components, no package manager, no graphical interface, and no
post-install configuration needed.

\section{What Makes Veilbox Different}

\subsection{Single-Binary Userspace}

BusyBox replaces over 300 standard Unix utilities (from coreutils, util-linux,
procps-ng, and many other packages) with a single statically-linked binary.
This reduces the root filesystem from hundreds of megabytes to approximately
90 megabytes (including containerd, runc, nerdctl (aliased as docker), and Dropbear binaries).

The BusyBox binary determines which utility to run based on argv{[}0{]}
(the name it was called with via symbolic links). For example, /bin/sh is a
symbolic link to /bin/busybox. When the kernel executes /bin/sh, it loads
/busybox and passes argv{[}0{]} = "/bin/sh"". BusyBox's main function
extracts the applet name (``sh'') from argv{[}0{]} and dispatches to the
shell implementation.

\subsection{Embedded Initramfs}

The root filesystem is compiled directly into the kernel binary via the
CONFIG\_INITRAMFS\_SOURCE kernel configuration option. This means the kernel
and its root filesystem are a single, atomic unit. There is no separate
initrd file to manage, no mkinitrd or dracut step, and no risk of the
initramfs getting out of sync with the kernel.

The initramfs is a compressed cpio archive that the kernel unpacks into a
tmpfs filesystem at boot time. The archive is generated by
gen\_initramfs\_list() in build.sh, which converts the rootfs directory
structure into the kernel's internal file list format.

\subsection{No Package Manager}

There is no rpm, dpkg, apk, or any other package manager. Software is
either compiled into the kernel, embedded in the initramfs at build time,
or run as a container via containerd. This dramatically reduces the attack
surface: there is no way for an attacker to install software, exploit a
package repository, or use package manager vulnerabilities.

The trade-off is that adding new software requires rebuilding the initramfs
and kernel. For a disposable VM or CI/CD worker, this is acceptable: changes
are made to build.sh and deployed by rebuilding the disk image.

\subsection{No Graphical Stack}

There is no X11, no Wayland, no framebuffer console, no display manager.
The system is managed entirely through two console interfaces: a serial
console (ttyS0) for headless access via QEMU's -nographic mode, and a
VGA text console (tty1) for monitor-and-keyboard access. This eliminates
hundreds of megabytes of graphics libraries, drivers, and utilities.

\subsection{Rootless Build}

The entire build process runs without sudo or root privileges. This is
achieved through several techniques:

\begin{itemize}[nosep]
  \item debugfs is used instead of mount to populate the ext4 filesystem.
        debugfs can write files into an ext4 image without mounting it,
        which avoids the need for loop device access (which requires root).
  \item The GRUB bootloader is installed via a custom Python script that
        writes boot.img and core.img directly to the disk image file,
        preserving the partition table written by sfdisk.
  \item Kernel source compilation runs as the regular user. No special
        permissions are needed for make or gcc.
  \item All downloaded binaries and tarballs are extracted without any
        special privileges.
\end{itemize}

The only exception is the dependency installation step (install\_deps()),
which may use sudo if running on Fedora with passwordless sudo configured.
On other distributions, or when sudo requires a password, the script prints
the package list and exits gracefully.

\section{Design Goals}

\subsection{No Installer}

A Veilbox user should never need to run an installer, answer configuration
questions, or wait for packages to download. The workflow is:

\begin{enumerate}[nosep]
  \item Clone the repository from GitHub.
  \item Run build.sh (approximately 5 minutes for a full build).
  \item Import the VDI into VirtualBox, or run test.sh for QEMU.
  \item Boot and log in.
\end{enumerate}

The entire system is pre-configured with sensible defaults: root password
(veiladmin), SSH host keys, containerd configuration, network settings,
and the bootloader are all generated during the build.

\subsection{Self-Contained}

The kernel, userspace, persistent storage, and bootloader all fit on a
single disk image. No external dependencies, no network requirements for
first boot, no cloud-init, no post-install scripts.

The disk image contains:
\begin{itemize}[nosep]
  \item MBR with GRUB bootloader (boot.img + core.img)
  \item ext4 partition with label VEILBOX
  \item /boot/vmlinuz (kernel with embedded initramfs)
  \item /boot/grub/grub.cfg (GRUB configuration)
  \item /boot/grub/*.mod (GRUB filesystem modules)
  \item /state/ directory for persistent data
\end{itemize}

\subsection{Minimal Attack Surface}

Every kernel feature and userspace utility that is not strictly necessary
is removed. The kernel has no sound, no DRM, no USB, no Wi-Fi, no Bluetooth,
no ATA or SCSI drivers. The defconfig baseline includes support for
approximately 10,000 kernel options; Veilbox reduces this aggressively.

In userspace, there are no compilers, no interpreters (other than the ash
shell), no setuid binaries, no graphical stack, and no network services
other than containerd and Dropbear. The initramfs is read-only (embedded
in the kernel), so runtime modification of system binaries is impossible.

\subsection{Container-Ready}

Containerd, runc, and nerdctl are pre-installed and configured out of the
box. The state partition is pre-created with directories for container
images, logs, and volumes. The first thing a user sees after logging in is
a working container host environment.

\subsection{Reproducible Build}

The build.sh script is designed to produce bit-identical results given the
same inputs. Version pinning ensures that every downloaded binary has a
fixed version number. The script includes a --clean flag that removes all
build artifacts and starts fresh. The initramfs file list is sorted for
deterministic ordering.

\section{Target Audience}

\begin{itemize}[nosep]
  \item \textbf{Self-hosters} looking for a lightweight test lab for
        containers. Veilbox boots in under 15 seconds and uses only
        2 GB of RAM, making it ideal for running on modest hardware.
  \item \textbf{DevOps engineers} needing minimal, disposable virtual
        machines for CI/CD pipelines. The health check mode (--check)
        integrates with automated testing.
  \item \textbf{Container enthusiasts} experimenting with containerd and
        nerdctl outside of the Docker ecosystem.
  \item \textbf{Embedded Linux students} studying how a kernel boots, how
        an initramfs works, how a bootloader loads an operating system,
        and how a userspace is assembled from scratch.
  \item \textbf{Security researchers} wanting a minimal, auditable Linux
        system where every included component serves a necessary purpose.
\end{itemize}

\section{What You Will Learn}

By reading this guide and building Veilbox, you will gain a deep
understanding of:

\begin{enumerate}[nosep]
  \item The Linux kernel build system: Kconfig, defconfig, menuconfig,
        merge\_config.sh, and the .config file format.
  \item The initramfs: cpio archives, the kernel file list format, device
        node creation, and how the kernel unpacks the root filesystem.
  \item The GRUB bootloader: boot.img, core.img, the blocklist format,
        the MBR partition table, and rootless installation.
  \item BusyBox: applet dispatch, the ash shell, busybox init, and how
        a single binary provides over 300 Unix utilities.
  \item Container runtimes: containerd architecture, gRPC API, runc OCI
        runtime, overlay filesystem layers, and nerdctl CLI.
  \item SSH with Dropbear: key generation, authentication methods, host
        key verification, and port forwarding.
  \item Linux namespaces and cgroups: PID, network, mount, user namespaces;
        CPU, memory, device cgroup controllers.
  \item QEMU virtualization: VirtIO paravirtualization, SLiRP networking,
        KVM acceleration, and disk image formats.
  \item Build automation: Shell scripting, error handling, caching, and
        depedency management for complex build pipelines.
\end{enumerate}
